\begin{abstract}
Iterated optimization-based bound tightening (OBBT) optimizes each
coordinate over a convex relaxation with an objective cutoff and rebuilds
it on the tightened box. For a valid relaxation family that is monotone
under box restriction, three things limit further tightening: the hull of
feasible points with objective at most the cutoff, boxes that the
relaxation cannot shrink, and the rate of approach to the limit. Near a
minimizer, the relaxation's second-order tangent model defines an
order-preserving map on box shapes. If its upper Collatz--Wielandt factor
is strictly below one, it bounds the local linear rate, and the limit has
width $O(\sqrt\epsilon)$ and leaves a relaxation gap $O(\epsilon)$ for
incumbent suboptimality $\epsilon$. Points of negative tangent-model value
on every face of a box shape certify stalling at every sufficiently small
scale. The model is exact for quadratic objectives with McCormick
relaxations, giving closed-form rates and stall conditions, including
stalls of strongly convex objectives under that relaxation; for factorable
functions with $C^2$ univariate factors, we prove a second-order expansion
of composite McCormick relaxations. Finite point pools feasible in the
relaxation rebuilt on their own hull certify boxes that no later round
removes, with exact cutoff thresholds; residual and sensitivity bounds
limit the remaining movement. In 120 prospective runs, an adaptive SCIP
policy using only current-round screening solved $20.9\%$ fewer auxiliary
LPs than a fixed policy but spent more propagator time; neither produced an
additional solve or demonstrated speedup over native SCIP. These
certificates bound a specified operator's movement, not search time.
\end{abstract}
